\documentclass[12pt]{article}
\usepackage{amsmath,amsfonts,fullpage,amsthm,graphicx}

\newtheorem*{theorem}{Theorem}
\newtheorem*{fact}{Fact}
\newtheorem*{goal}{Goal}
\newtheorem*{idea}{Idea}
\newtheorem*{defn}{Definition}
\newtheorem*{ex}{Example}

%Style section
%\setlength{\textheight}{10in}
%\setlength{\textwidth}{7in}
%\hoffset=-0.75in
 \pagestyle{plain}
% \setlength{\topmargin}{0in}
% \setlength{\headsep}{0in}
% \setlength{\oddsidemargin}{0in}
% \setlength{\evensidemargin}{0in}

\begin{document}


\pagestyle{empty} %use this to get rid of page numbers

\noindent
{Math 166 \hskip 1.5 truein  Names:\ \hrulefill}

\noindent
%\vskip 0.3truein
%\bigskip
%\noindent
{A study of series}

\noindent \hskip 3.5 truein  Section Number:\ \hrulefill
\medskip

\begin{goal} \rm
Explore methods of determining whether series converge or diverge. 
\end{goal}
%\smallskip

\begin{enumerate}
\setcounter{enumi}{-1}
%Name all the tests you can think of for convergence/divergence of series in our series toolbox. Put a * by the tests that only work for series with positive terms.
\iffalse
\item
\begin{itemize}
\item A \textbf{sequence} is a~ \rule{.8in}{.4pt} ~of numbers.  A sequence converges if its~ \rule{.8in}{.4pt} ~have a limit.
\item A \textbf{series} is a~ \rule{.6in}{.4pt} ~of numbers. A series converges if its~ \rule{.8in}{.4pt} \rule{.2in}{.4pt} \rule{.8in}{.4pt} \rule{.5in}{.4pt} ~has a limit.
\end{itemize}
\fi

\item State the direct comparison  in your own words. What conditions do you need to check in order to use this test?

%\vspace{.5in}
%\item In your team, each person should attempt to use \textbf{one} of the following tests to determine whether the series  converges or diverges: $n$th term divergence test, Comparison test, Limit comparison test, Integral test. Which test(s) worked well for this series? 



\vfill
You can also use the following test.

\bigskip
\noindent
\underline{\textbf{Limit comparison test}}\\

Suppose $\displaystyle\sum_{n=1}^\infty a_n$   and $\displaystyle\sum_{n=1}^\infty b_n$ are series with positive terms ($a_n\geq 0$ and $b_n\geq 0$) such that $\displaystyle\lim_{n\to \infty}\frac{a_n}{b_n} =L$  for some number $L$. 
\begin{enumerate}
    \item If $0<L<\infty$, then the series either both converge or  both diverge. 
    \item If $L = \infty$ and $\displaystyle\sum_{n=1}^\infty a_n$  converges, then $\displaystyle\sum_{n=1}^\infty b_n$ converges.
    \item If $L = 0$ and $\displaystyle\sum_{n=1}^\infty b_n$ converges, then $\displaystyle\sum_{n=1}^\infty a_n$ converges. 
\end{enumerate}

\pagebreak
\item Use any of the following tests to determine whether the following series converge or diverge: $n$th term divergence test, Integral test, Comparison test, Limit comparison test. Be sure to verify the hypotheses, state your conclusion, and say which test you used. 
%You should approach each problem together as a team, but you may each try different tests to more quickly determine what test will work in each problem.
\begin{enumerate}
%\item $\displaystyle\sum_{n=1}^{\infty} (1.01)^n$

%\vspace{3.5in}
%\pagebreak
%\vspace{1.5in}

%\item $\displaystyle\sum_{n=1}^{\infty} e^{-3n}$

%\vfill
\item $\displaystyle\sum_{n=2}^{\infty} \frac{1}{\sqrt{n^4+n^3}}$

\vfill

%\pagebreak
%\item $\displaystyle\sum_{n=1}^{\infty} \frac{7^n+1}{2^n}$

%\vfill
%\item $\displaystyle\sum_{n=2}^{\infty} \frac{n}{n^2+1}$

%\vfill
%\item $\displaystyle\sum_{n=1}^{\infty} \frac{\cos(n)}{n^2}$ (Hint: Can you use the alternating series test? Why not? Instead, check for absolute convergence.)
\item $\displaystyle\sum_{n=1}^{\infty} \frac{1}{4+5^n}$

\vfill
\item 
$\displaystyle\sum_{n=1}^{\infty} \frac{n^3-5n+1}{7n^5+2}$

\vfill
\pagebreak

\item $\displaystyle\sum_{n=1}^{\infty} \frac{e^{-n}}{n+\cos^2(n)}$
\vfill

\item $\displaystyle\sum_{n=2}^{\infty} \frac{4n^2+n}{\sqrt[3]{n^7+n^3}}$
\vfill 

\item $\displaystyle\sum_{k=2}^{\infty} \sqrt{\frac{k}{k^3+1}}$
\vfill
\end{enumerate} 

\end{enumerate}

\end{document}
